\section{Related Work}
\label{sec:related}

We organize related work around four lines most relevant to our problem setting: multi-person pose estimation, transformer-based context modeling, multi-scale small-object representation, and sports-field localization with localization-oriented keypoint refinement. This organization is intentional. Our method is not a generic pose estimator, but a single-stage localization system designed specifically to reduce occluded players, localization error, and false positive players in crowded soccer scenes.

\subsection{Multi-Person Pose Estimation}
\label{subsec:pose}

Multi-person pose estimation methods are commonly divided into top-down, bottom-up, and one-stage paradigms. Top-down methods first detect person boxes and then estimate keypoints for each cropped instance. Representative models such as SimpleBaseline~\cite{xiao2018simple}, HRNet~\cite{sun2019deep}, ViTPose~\cite{xu2022vitpose}, and RTMPose~\cite{jiang2023rtmpose} achieve strong per-instance accuracy by dedicating substantial computation to each detected person. However, their computational cost grows with the number of players, which is undesirable in broadcast soccer scenes where many athletes may appear simultaneously and crowding is common~\cite{fang2017rmpe}.

Bottom-up methods reverse this order by first predicting all keypoints and then grouping them into person instances. OpenPose~\cite{cao2017realtime}, Associative Embedding~\cite{newell2017associative}, HigherHRNet~\cite{cheng2020higherhrnet}, and DEKR~\cite{geng2021bottom} show that this paradigm scales better with person count. Yet the grouping stage itself becomes fragile in dense interactions, where neighboring players overlap heavily and association ambiguity increases. For crowded soccer imagery, this means that bottom-up designs alleviate one bottleneck while introducing another.

One-stage methods combine detection and pose prediction in a unified forward pass. CenterNet~\cite{zhou2019objects}, DirectPose~\cite{tian2019directpose}, and YOLO-Pose~\cite{maji2022yolo} illustrate the efficiency of this route. More recently, the Ultralytics family has strengthened one-stage pose estimation through YOLOv8-Pose~\cite{jocher2023yolov8}, YOLO11-Pose~\cite{jocher2024yolo11}, and YOLO26-Pose~\cite{ultralytics2025yolo26}. We build on YOLO26-Pose because it provides the strongest real-time one-stage baseline in our setting. Our contribution is therefore not to replace the one-stage paradigm, but to adapt it for world-coordinate soccer localization by explicitly targeting three failure categories that standard pose pipelines do not directly optimize for.

\subsection{Transformer-Based Detection and Context Modeling}
\label{subsec:transformer}

Transformer-based detectors provide a natural reference for modeling long-range context in crowded scenes. DETR~\cite{carion2020end} demonstrates that set prediction with global attention can remove hand-crafted detection components such as anchors and non-maximum suppression. Subsequent variants improve efficiency and convergence: Deformable DETR~\cite{zhu2021deformable} introduces sparse multi-scale deformable attention, DAB-DETR~\cite{liu2022dabdetr} adds dynamic anchor queries, DN-DETR~\cite{li2022dn} accelerates training with denoising, DINO~\cite{zhang2023dino} further strengthens optimization and accuracy, and Co-DETR~\cite{zong2023detrs} enriches encoder supervision through collaborative assignments.

Among these methods, RT-DETR~\cite{zhao2024detrs} is most relevant to our design because it shows how transformer-style reasoning can be injected efficiently into a real-time detector. Its efficient hybrid encoder separates intra-scale and cross-scale interaction and uses AIFI to model long-range dependencies on the highest-level feature map. This design is conceptually close to our use of AIFI at P5. However, our goal differs from generic object detection. We do not replace the whole one-stage pose backbone with a transformer detector. Instead, we retain YOLO26x-Pose and introduce transformer-style modeling only at the stages that directly address our task: global semantic disambiguation for false positives, cross-scale enhancement for difficult players, and residual refinement for localization error.

\subsection{Multi-Scale Feature Fusion and Small-Object Representation}
\label{subsec:multiscale}

Multi-scale fusion is central to detecting players across large scale variation. FPN~\cite{lin2017feature} establishes the standard top-down pyramid with lateral connections, while PANet~\cite{liu2018path} strengthens information flow through a bottom-up path. BiFPN~\cite{tan2020efficientdet} and NAS-FPN~\cite{ghiasi2019nasfpn} further improve learned feature aggregation. These designs have become standard components in modern YOLO-style detectors and remain strong generic solutions for multi-scale recognition.

The limitation in our setting is not the absence of a feature pyramid, but the way weak small-instance evidence is propagated through it. When athletes are far from the camera, their P3 responses are fragile, and fixed convolutional fusion may not preserve enough semantically meaningful support from coarser levels. This weakness becomes especially visible in crowded soccer scenes, where distant players are both small and partially entangled with nearby instances. Our cross-scale deformable attention module addresses this gap by allowing P3 queries to sample sparse, informative support from P3, P4, and P5 jointly, following the sparse-sampling intuition of Deformable DETR~\cite{zhu2021deformable} but adapting it to one-stage pose localization.

\subsection{Keypoint Refinement and Sports-Field Localization}
\label{subsec:sports}

A second relevant line of work concerns refinement after coarse prediction. Integral pose regression~\cite{sun2018integral}, DarkPose~\cite{zhang2020distribution}, and UDP~\cite{huang2020devil} show that substantial gains can come from correcting coordinate-level bias rather than only improving coarse heatmaps. Cascaded and refinement-based designs such as CPN~\cite{chen2018cascaded}, RSN~\cite{cai2020learning}, and PoseFix~\cite{moon2019posefix} further demonstrate that post-hoc correction is particularly useful when the base predictor is already strong. Our local refiner follows the same coarse-to-fine philosophy, but with a different task target: instead of improving general human pose quality, it specifically corrects the \texttt{pelvis} and \texttt{pelvis\_ground} coordinates whose projection determines LocSim.

The task setting itself is rooted in sports-field localization and player positioning. SoccerNet~\cite{giancola2018soccernet}, SoccerNet-v2~\cite{deliege2021soccernetv2}, SoccerNet-v3~\cite{cioppa2022scaling}, and the Game State Reconstruction benchmark~\cite{somers2024soccernet} establish the importance of athlete localization, calibration, and field-level reasoning in soccer analytics. Related calibration and field-registration approaches include deep structured models~\cite{homayounfar2017sports}, optimization through learned errors~\cite{jiang2020optimizing}, end-to-end calibration~\cite{sha2020end2end}, synthetic-data-driven calibration~\cite{chen2019sports}, robust field registration~\cite{nie2021robust}, and TVCalib~\cite{theiner2023tvcalib}. Closest to our setting is SSKIT~\cite{ardo2025}, which introduces the SynLoc task of single-frame world-coordinate athlete detection and localization on synthetic 4K soccer imagery with LocSim as the evaluation metric. CLEAR is designed precisely for this setting: it keeps the single-frame one-stage formulation of SynLoc, but strengthens the model where SynLoc is most demanding, namely under occlusion, projection-sensitive localization error, and false positive clutter.
